\section{Discussion and Limitations}

\subsection{Interventions as conditional measurements}

Raw displacement is insufficient within residual attenuation, and functional
damage is more informative, yet joint KL/NLL matching leaves a stable
block/noise difference. Intervention identity therefore retains information
not captured by either severity description; it is not itself an explanation.
A layer may look substitutable under one procedure and not another because the
procedures trace different paths through activation space. Intervention claims
should consequently report the model, data, family, magnitude convention,
functional-damage budget, response metric, and support diagnostics rather than
describe protocol-independent ``layer redundancy.''

\subsection{What remains unidentified}

The two families induce strongly different hidden perturbation directions, but
frozen observational direction features do not pass their preregistered
explanatory gates. A planned counterfactual direction test fails its support
gate before top-1 outcome reveal. We therefore make no causal claim about
internal direction; the blinded feasibility results remain in
Appendix~\ref{app:direction}.

\subsection{Limitations}

\paragraph{Scope.}
Evidence is limited to 70M/160M \model{} replicas, two English evaluation
streams with one tokenizer, and two intervention families; cross-family damage
matching is 160M and WikiText-2 only. Corpus magnitudes differ by about 23\%,
and replacement, interchange, attention-only, MLP-only, editing, other
architectures, and billion-parameter models remain untested.

\paragraph{Construct validity and heterogeneity.}
$S$ is discontinuous teacher-forced top-1 agreement, not semantic equivalence,
downstream capability, or pruning safety. KL, NLL, geometry, and $D_S$ share
logits, so matched associations may retain common metric geometry. Although the
aggregate family residual is seed-stable, 24/103 cells reverse sign and layer
effects vary; low support at layer 7 limits location-specific interpretation.

\paragraph{Mechanism and assay repair.}
The frozen tests weaken raw-magnitude, family-invariant damage, and simple
output-geometry accounts but do not identify an internal mechanism. The first
prospective reveal was invalidated by inconsistent FP16 top-1 tie handling; we
stopped, withdrew it, froze a family-independent lowest-index rule, and
recomputed without changing matches or exclusions. The corrected conclusion
survives, but the 12.49\% point-estimate change is material; full provenance is
in Appendix~\ref{app:repair}.
